\documentclass[11pt,a4paper]{article}
\usepackage[utf8]{inputenc}
\usepackage[T1]{fontenc}
\usepackage[spanish,es-minimal,es-noshorthands]{babel}
\AtBeginDocument{\renewcommand{\sin}{\operatorname{sen}}\renewcommand{\arcsin}{\operatorname{arcsen}}}
\usepackage{lmodern}
\usepackage[sfdefault,lining]{FiraSans}
\usepackage[a4paper,top=2.3cm,bottom=2.2cm,left=1.8cm,right=1.8cm,headheight=1.3cm,headsep=0.4cm,footskip=1cm]{geometry}
\usepackage{amsmath,amssymb}
\usepackage{icomma}
\usepackage{graphicx}
\usepackage[table]{xcolor}
\usepackage{tikz}
\usepackage{fontawesome5}
\usepackage{ccicons}
\usepackage{fancyhdr}
\usepackage{lastpage}
\usepackage{tasks}
\usepackage{needspace}
\usepackage[most]{tcolorbox}
\usepackage{array}
\usepackage[hidelinks]{hyperref}
\AtBeginDocument{\hypersetup{pdftitle={Ficha de ejercicios de la unidad 7: Derivadas y aplicaciones. 1º Bach.},pdfsubject={Matemáticas CCSS 1º Bach.},pdfkeywords={tasa de variación media, derivada, función derivada, reglas de derivación, regla de la cadena, recta tangente, recta normal, monotonía, máximos y mínimos, representación gráfica, optimización}}}

\definecolor{navy}{HTML}{1B2748}
\definecolor{accent}{HTML}{E8603A}
\definecolor{teal}{HTML}{1F9E90}
\definecolor{soft}{HTML}{EEF1F8}
\definecolor{muted}{HTML}{5D6685}

\setlength{\parindent}{0pt}
\setlength{\parskip}{3pt}
\settasks{label=\alph*), label-format=\bfseries\color{navy}, label-width=1.3em, item-indent=1.8em, column-sep=0.6em, after-item-skip=5pt, before-skip=3pt, after-skip=4pt}

% ---------- capçalera i peu ----------
\newcommand{\llicencia}{\href{https://creativecommons.org/licenses/by-nc-sa/4.0/deed.es}{\ccbyncsa}\ \ Licencia Creative Commons Atribución-NoComercial-CompartirIgual 4.0 Internacional (CC BY-NC-SA 4.0)}
\pagestyle{fancy}
\fancyhf{}
\fancyhead[L]{\raisebox{-0.35\height}{\includegraphics[height=0.95cm]{logo-gva.png}}}
\fancyhead[C]{\small\color{muted}Ficha de ejercicios · Unidad 7: Derivadas y aplicaciones · 1\textsuperscript{o} Bach.}
\fancyhead[R]{\raisebox{-0.35\height}{\includegraphics[height=0.95cm]{logo-vermellar.png}}}
\fancyfoot[L]{\scriptsize\color{muted}\llicencia}
\fancyfoot[R]{\small\color{muted}\thepage}
\renewcommand{\headrulewidth}{0.4pt}
\renewcommand{\headrule}{\color{muted!40}\hrule height \headrulewidth}
\fancypagestyle{primera}{\fancyhf{}\renewcommand{\headrulewidth}{0pt}
  \fancyfoot[L]{\scriptsize\color{muted}\llicencia}
  \fancyfoot[R]{\small\color{muted}\thepage}}

% ---------- nivells i exercicis ----------
\newcommand{\estrelles}[1]{{\color{accent}\ifcase#1\or\faStar\faStar[regular]\faStar[regular]\or\faStar\faStar\faStar[regular]\or\faStar\faStar\faStar\fi}}
\newcounter{exer}
\newcommand{\ex}[1]{\par\needspace{4\baselineskip}\medskip\refstepcounter{exer}\noindent{\color{navy}\bfseries\theexer.}\ \ }

\newcommand{\seccio}[2]{\par\needspace{16\baselineskip}\bigskip
  \begin{tcolorbox}[colback=navy,colframe=navy,coltext=white,arc=2mm,boxrule=0pt,left=3mm,right=3mm,top=1.2mm,bottom=1.2mm]
  {\large\bfseries #1\quad #2}
  \end{tcolorbox}\vspace{-2pt}}

\newcommand{\exf}[1]{\par\medskip\refstepcounter{exer}\noindent{\color{navy}\bfseries\theexer.}\ \ }
\providecommand{\gt}{>}\providecommand{\lt}{<}
% quadre de recordatori (només a la fitxa d'exercicis)
\newcommand{\recorda}[1]{\par\smallskip\begin{tcolorbox}[colback=accent!7,colframe=accent!55,arc=2mm,boxrule=0.5pt,left=2.5mm,right=2.5mm,top=1mm,bottom=1mm]\small{\color{accent}\faLightbulb[regular]}\ \textbf{Recuerda:}\ #1\end{tcolorbox}}
\newcommand{\res}[1]{#1}\newcommand{\rest}[1]{#1}

\def\fmneed{11}% que el títol d'una secció no quede sol al final del full
% =====================================================================
% Capçalera FentMates de les fitxes: un color per a cada curs.
% Els colors són els de la pàgina inicial (LEVELS a assets/js/materials.js);
% el logo ve de _eines/logo/pdf/ (python3 _eines/logo/genera.py).
%
% Ús, al final del preàmbul de la unitat:
%   % =====================================================================
% Capçalera FentMates de les fitxes: un color per a cada curs.
% Els colors són els de la pàgina inicial (LEVELS a assets/js/materials.js);
% el logo ve de _eines/logo/pdf/ (python3 _eines/logo/genera.py).
%
% Ús, al final del preàmbul de la unitat:
%   % =====================================================================
% Capçalera FentMates de les fitxes: un color per a cada curs.
% Els colors són els de la pàgina inicial (LEVELS a assets/js/materials.js);
% el logo ve de _eines/logo/pdf/ (python3 _eines/logo/genera.py).
%
% Ús, al final del preàmbul de la unitat:
%   \input{../comu/fentmates.tex}
%   \fmcurs{1eso}
% i, al començament del document:
%   \capcalera{Matemàtiques}{1r ESO}{Fitxa d'exercicis · Unitat 2}{Divisibilitat}
% Opcional abans de \input: \def\fmneed{14} (espai mínim abans de cada secció).
% =====================================================================
\makeatletter
\def\fm@logodir{../../logo/pdf}
\providecommand{\fmneed}{10}
\def\fm@cursnom{}

% clau del curs, color i color del text damunt del color
\newcommand{\fm@curs}[3]{\@namedef{fm@color@#1}{#2}\@namedef{fm@text@#1}{#3}}
\fm@curs{1eso}{E6B84A}{navy}
\fm@curs{2eso}{3B82F6}{white}
\fm@curs{3eso}{FF7A50}{white}
\fm@curs{4eso-a}{A78BFA}{white}
\fm@curs{4eso-b}{6D28D9}{white}
\fm@curs{1batx-mi}{9B1C1C}{white}
\fm@curs{1batx-ccss}{EF6B6B}{white}
\fm@curs{2batx-mii}{1E6B47}{white}
\fm@curs{2batx-ccss}{2BB3A3}{white}

\newcommand{\fmcurs}[1]{%
  \@ifundefined{fm@color@#1}{\PackageError{fentmates}{Curs desconegut: #1}{}}{}%
  \edef\fm@tmp{\noexpand\definecolor{curs}{HTML}{\@nameuse{fm@color@#1}}%
               \noexpand\colorlet{curstext}{\@nameuse{fm@text@#1}}}\fm@tmp
  \colorlet{cursfosc}{curs!72!black}%
  \def\fm@logo{\fm@logodir/vertical-#1.pdf}%
  \def\fm@marca{\fm@logodir/marca-#1.pdf}}

% ---------- capçalera de la primera pàgina ----------
% #1 matèria (petit), #2 curs, #3 tipus de document i unitat, #4 títol
\newcommand{\capcalera}[4]{%
  \gdef\fm@cursnom{#2}%
  \par\noindent
  \begin{tcolorbox}[colback=curs,colframe=curs,boxrule=0pt,arc=2.2mm,width=3.5cm,height=2.15cm,
      valign=center,left=3mm,right=2mm,top=1mm,bottom=1mm,nobeforeafter,box align=center]
    \color{curstext}{\footnotesize #1}\par\vspace{1pt}{\fontsize{19}{21}\selectfont\bfseries #2\par}%
  \end{tcolorbox}\hspace{5mm}%
  \begin{minipage}[c]{\dimexpr\linewidth-3.5cm-5mm-2.9cm\relax}
    {\color{cursfosc}\small\bfseries\MakeUppercase{#3}}\par\vspace{2pt}
    {\color{navy}\fontsize{20}{23}\selectfont\bfseries #4\par}
  \end{minipage}\hfill
  \begin{minipage}[c]{2.6cm}\raggedleft\includegraphics[width=2.6cm]{\fm@logo}\end{minipage}\par
  \vspace{2pt}}

% ---------- barres de secció del color del curs ----------
\renewcommand{\seccio}[2]{\par\needspace{\fmneed\baselineskip}\bigskip
  \begin{tcolorbox}[colback=curs,colframe=curs,coltext=curstext,arc=2mm,boxrule=0pt,left=3mm,right=3mm,top=1.2mm,bottom=1.2mm]
  {\large\bfseries #1\quad #2}
  \end{tcolorbox}\vspace{-2pt}}

% ---------- estrelles de dificultat del color del curs (a Batxillerat no se'n posen) ----------
\@ifundefined{estrelles}{}{%
  \renewcommand{\estrelles}[1]{{\color{curs}\ifcase#1\or\faStar\faStar[regular]\faStar[regular]\or\faStar\faStar\faStar[regular]\or\faStar\faStar\faStar\fi}}}

% ---------- capçalera i peu de les pàgines ----------
% primera pàgina: logos institucionals al peu
\fancypagestyle{primera}{\fancyhf{}\renewcommand{\headrulewidth}{0pt}
  \fancyfoot[L]{\raisebox{-0.32\height}{\includegraphics[height=0.78cm]{logo-gva.png}}\hspace{5mm}%
    \raisebox{-0.32\height}{\includegraphics[height=0.78cm]{logo-vermellar.png}}}
  \fancyfoot[R]{\scriptsize\color{muted}\href{https://creativecommons.org/licenses/by-nc-sa/4.0/deed.ca}{\ccbyncsa}\ \ CC BY-NC-SA 4.0 · fentmates.org\qquad\small\thepage}}
% la resta: el símbol a l'esquerra i el curs a la dreta, amb el títol al centre com abans
\fancyhead[L]{\raisebox{-0.3\height}{\includegraphics[height=0.62cm]{\fm@marca}}}
\fancyhead[R]{\small\bfseries\color{cursfosc}\fm@cursnom}
\renewcommand{\headrule}{\color{curs!45}\hrule height \headrulewidth}
\makeatother

%   \fmcurs{1eso}
% i, al començament del document:
%   \capcalera{Matemàtiques}{1r ESO}{Fitxa d'exercicis · Unitat 2}{Divisibilitat}
% Opcional abans de \input: \def\fmneed{14} (espai mínim abans de cada secció).
% =====================================================================
\makeatletter
\def\fm@logodir{../../logo/pdf}
\providecommand{\fmneed}{10}
\def\fm@cursnom{}

% clau del curs, color i color del text damunt del color
\newcommand{\fm@curs}[3]{\@namedef{fm@color@#1}{#2}\@namedef{fm@text@#1}{#3}}
\fm@curs{1eso}{E6B84A}{navy}
\fm@curs{2eso}{3B82F6}{white}
\fm@curs{3eso}{FF7A50}{white}
\fm@curs{4eso-a}{A78BFA}{white}
\fm@curs{4eso-b}{6D28D9}{white}
\fm@curs{1batx-mi}{9B1C1C}{white}
\fm@curs{1batx-ccss}{EF6B6B}{white}
\fm@curs{2batx-mii}{1E6B47}{white}
\fm@curs{2batx-ccss}{2BB3A3}{white}

\newcommand{\fmcurs}[1]{%
  \@ifundefined{fm@color@#1}{\PackageError{fentmates}{Curs desconegut: #1}{}}{}%
  \edef\fm@tmp{\noexpand\definecolor{curs}{HTML}{\@nameuse{fm@color@#1}}%
               \noexpand\colorlet{curstext}{\@nameuse{fm@text@#1}}}\fm@tmp
  \colorlet{cursfosc}{curs!72!black}%
  \def\fm@logo{\fm@logodir/vertical-#1.pdf}%
  \def\fm@marca{\fm@logodir/marca-#1.pdf}}

% ---------- capçalera de la primera pàgina ----------
% #1 matèria (petit), #2 curs, #3 tipus de document i unitat, #4 títol
\newcommand{\capcalera}[4]{%
  \gdef\fm@cursnom{#2}%
  \par\noindent
  \begin{tcolorbox}[colback=curs,colframe=curs,boxrule=0pt,arc=2.2mm,width=3.5cm,height=2.15cm,
      valign=center,left=3mm,right=2mm,top=1mm,bottom=1mm,nobeforeafter,box align=center]
    \color{curstext}{\footnotesize #1}\par\vspace{1pt}{\fontsize{19}{21}\selectfont\bfseries #2\par}%
  \end{tcolorbox}\hspace{5mm}%
  \begin{minipage}[c]{\dimexpr\linewidth-3.5cm-5mm-2.9cm\relax}
    {\color{cursfosc}\small\bfseries\MakeUppercase{#3}}\par\vspace{2pt}
    {\color{navy}\fontsize{20}{23}\selectfont\bfseries #4\par}
  \end{minipage}\hfill
  \begin{minipage}[c]{2.6cm}\raggedleft\includegraphics[width=2.6cm]{\fm@logo}\end{minipage}\par
  \vspace{2pt}}

% ---------- barres de secció del color del curs ----------
\renewcommand{\seccio}[2]{\par\needspace{\fmneed\baselineskip}\bigskip
  \begin{tcolorbox}[colback=curs,colframe=curs,coltext=curstext,arc=2mm,boxrule=0pt,left=3mm,right=3mm,top=1.2mm,bottom=1.2mm]
  {\large\bfseries #1\quad #2}
  \end{tcolorbox}\vspace{-2pt}}

% ---------- estrelles de dificultat del color del curs (a Batxillerat no se'n posen) ----------
\@ifundefined{estrelles}{}{%
  \renewcommand{\estrelles}[1]{{\color{curs}\ifcase#1\or\faStar\faStar[regular]\faStar[regular]\or\faStar\faStar\faStar[regular]\or\faStar\faStar\faStar\fi}}}

% ---------- capçalera i peu de les pàgines ----------
% primera pàgina: logos institucionals al peu
\fancypagestyle{primera}{\fancyhf{}\renewcommand{\headrulewidth}{0pt}
  \fancyfoot[L]{\raisebox{-0.32\height}{\includegraphics[height=0.78cm]{logo-gva.png}}\hspace{5mm}%
    \raisebox{-0.32\height}{\includegraphics[height=0.78cm]{logo-vermellar.png}}}
  \fancyfoot[R]{\scriptsize\color{muted}\href{https://creativecommons.org/licenses/by-nc-sa/4.0/deed.ca}{\ccbyncsa}\ \ CC BY-NC-SA 4.0 · fentmates.org\qquad\small\thepage}}
% la resta: el símbol a l'esquerra i el curs a la dreta, amb el títol al centre com abans
\fancyhead[L]{\raisebox{-0.3\height}{\includegraphics[height=0.62cm]{\fm@marca}}}
\fancyhead[R]{\small\bfseries\color{cursfosc}\fm@cursnom}
\renewcommand{\headrule}{\color{curs!45}\hrule height \headrulewidth}
\makeatother

%   \fmcurs{1eso}
% i, al començament del document:
%   \capcalera{Matemàtiques}{1r ESO}{Fitxa d'exercicis · Unitat 2}{Divisibilitat}
% Opcional abans de \input: \def\fmneed{14} (espai mínim abans de cada secció).
% =====================================================================
\makeatletter
\def\fm@logodir{../../logo/pdf}
\providecommand{\fmneed}{10}
\def\fm@cursnom{}

% clau del curs, color i color del text damunt del color
\newcommand{\fm@curs}[3]{\@namedef{fm@color@#1}{#2}\@namedef{fm@text@#1}{#3}}
\fm@curs{1eso}{E6B84A}{navy}
\fm@curs{2eso}{3B82F6}{white}
\fm@curs{3eso}{FF7A50}{white}
\fm@curs{4eso-a}{A78BFA}{white}
\fm@curs{4eso-b}{6D28D9}{white}
\fm@curs{1batx-mi}{9B1C1C}{white}
\fm@curs{1batx-ccss}{EF6B6B}{white}
\fm@curs{2batx-mii}{1E6B47}{white}
\fm@curs{2batx-ccss}{2BB3A3}{white}

\newcommand{\fmcurs}[1]{%
  \@ifundefined{fm@color@#1}{\PackageError{fentmates}{Curs desconegut: #1}{}}{}%
  \edef\fm@tmp{\noexpand\definecolor{curs}{HTML}{\@nameuse{fm@color@#1}}%
               \noexpand\colorlet{curstext}{\@nameuse{fm@text@#1}}}\fm@tmp
  \colorlet{cursfosc}{curs!72!black}%
  \def\fm@logo{\fm@logodir/vertical-#1.pdf}%
  \def\fm@marca{\fm@logodir/marca-#1.pdf}}

% ---------- capçalera de la primera pàgina ----------
% #1 matèria (petit), #2 curs, #3 tipus de document i unitat, #4 títol
\newcommand{\capcalera}[4]{%
  \gdef\fm@cursnom{#2}%
  \par\noindent
  \begin{tcolorbox}[colback=curs,colframe=curs,boxrule=0pt,arc=2.2mm,width=3.5cm,height=2.15cm,
      valign=center,left=3mm,right=2mm,top=1mm,bottom=1mm,nobeforeafter,box align=center]
    \color{curstext}{\footnotesize #1}\par\vspace{1pt}{\fontsize{19}{21}\selectfont\bfseries #2\par}%
  \end{tcolorbox}\hspace{5mm}%
  \begin{minipage}[c]{\dimexpr\linewidth-3.5cm-5mm-2.9cm\relax}
    {\color{cursfosc}\small\bfseries\MakeUppercase{#3}}\par\vspace{2pt}
    {\color{navy}\fontsize{20}{23}\selectfont\bfseries #4\par}
  \end{minipage}\hfill
  \begin{minipage}[c]{2.6cm}\raggedleft\includegraphics[width=2.6cm]{\fm@logo}\end{minipage}\par
  \vspace{2pt}}

% ---------- barres de secció del color del curs ----------
\renewcommand{\seccio}[2]{\par\needspace{\fmneed\baselineskip}\bigskip
  \begin{tcolorbox}[colback=curs,colframe=curs,coltext=curstext,arc=2mm,boxrule=0pt,left=3mm,right=3mm,top=1.2mm,bottom=1.2mm]
  {\large\bfseries #1\quad #2}
  \end{tcolorbox}\vspace{-2pt}}

% ---------- estrelles de dificultat del color del curs (a Batxillerat no se'n posen) ----------
\@ifundefined{estrelles}{}{%
  \renewcommand{\estrelles}[1]{{\color{curs}\ifcase#1\or\faStar\faStar[regular]\faStar[regular]\or\faStar\faStar\faStar[regular]\or\faStar\faStar\faStar\fi}}}

% ---------- capçalera i peu de les pàgines ----------
% primera pàgina: logos institucionals al peu
\fancypagestyle{primera}{\fancyhf{}\renewcommand{\headrulewidth}{0pt}
  \fancyfoot[L]{\raisebox{-0.32\height}{\includegraphics[height=0.78cm]{logo-gva.png}}\hspace{5mm}%
    \raisebox{-0.32\height}{\includegraphics[height=0.78cm]{logo-vermellar.png}}}
  \fancyfoot[R]{\scriptsize\color{muted}\href{https://creativecommons.org/licenses/by-nc-sa/4.0/deed.ca}{\ccbyncsa}\ \ CC BY-NC-SA 4.0 · fentmates.org\qquad\small\thepage}}
% la resta: el símbol a l'esquerra i el curs a la dreta, amb el títol al centre com abans
\fancyhead[L]{\raisebox{-0.3\height}{\includegraphics[height=0.62cm]{\fm@marca}}}
\fancyhead[R]{\small\bfseries\color{cursfosc}\fm@cursnom}
\renewcommand{\headrule}{\color{curs!45}\hrule height \headrulewidth}
\makeatother

\fmcurs{1batx-ccss}

\begin{document}
\thispagestyle{primera}

\capcalera{Matemáticas CCSS}{1\textsuperscript{o} Bach.}{Ficha de ejercicios · Unidad 7}{\raggedright Derivadas y aplicaciones}
\vspace{10pt}
Nombre y apellidos: \rule{6.9cm}{0.4pt}\hfill Grupo: \rule{1.6cm}{0.4pt}\hfill Fecha: \rule{2.2cm}{0.4pt}
\vspace{8pt}

\begin{tcolorbox}[colback=soft,colframe=soft,arc=2mm,left=3mm,right=3mm,top=1.5mm,bottom=1.5mm]
\small Resuelve los ejercicios en el cuaderno mostrando todos los pasos y simplifica siempre las derivadas. Las secciones siguen el orden de los apuntes. Puedes comprobar las derivadas con la calculadora de derivadas de los apuntes interactivos.
\end{tcolorbox}

% =====================================================
\seccio{1}{Tasa de variación media}

\exf{1} Halla la TVM de la función $y=x^2-8x+12$ en los intervalos $[1,2]$, $[1,5]$ y $[1,8]$. Después, hállala en el intervalo variable $[1,1+h]$ y comprueba que, dando a $h$ los valores 1, 4 y 7, obtienes los resultados anteriores.\recorda{$\text{TVM}[a,b]=\dfrac{f(b)-f(a)}{b-a}$ \ es la pendiente del segmento que une los puntos $(a,f(a))$ y $(b,f(b))$. Si la TVM es positiva, la función ha crecido de media en el intervalo.}

\ex{1} Lanzamos una pelota hacia arriba y su altura (en metros) después de $t$ segundos es $s(t)=20t-5t^2$. Calcula la velocidad media en los intervalos $[0,1]$, $[1,2]$ y $[2,4]$. ¿Qué indica el signo del último resultado?

\ex{1} Calcula la TVM de cada función en el intervalo indicado:\begin{tasks}(3)
  \task $f(x)=\sqrt x$ \ en $[1,4]$
  \task $g(x)=2^x$ \ en $[0,3]$
  \task $h(x)=\dfrac1x$ \ en $[1,2]$
\end{tasks}


% =====================================================
\seccio{2}{Derivada en un punto y función derivada}

\exf{1} Calcula, aplicando la definición, las derivadas siguientes:\begin{tasks}(1)
  \task La derivada de $y=\dfrac1x$ en el punto $x=-2$.
  \task La derivada de $y=-2x+4$ en los puntos de abscisas $-3$, $0$, $4$ y $7$. ¿Por qué sale siempre lo mismo?
  \task La derivada de $y=3x^2-5x+1$ en los puntos de abscisas $-2$, $1$ y $0$.
\end{tasks}\recorda{$f'(a)=\displaystyle\lim_{h\to0}\frac{f(a+h)-f(a)}{h}$: es la pendiente de la recta tangente a la gráfica en el punto $(a,f(a))$.}

\ex{2} Halla la función derivada aplicando la definición $f'(x)=\displaystyle\lim_{h\to0}\frac{f(x+h)-f(x)}{h}$:\begin{tasks}(3)
  \task $f(x)=x^2+3x$
  \task $f(x)=\dfrac{2}{x+1}$
  \task $f(x)=\sqrt{x+1}$
\end{tasks}

\ex{1} Calcula por definición la pendiente de la recta tangente a $y=x^2$ en $x=3$ y escribe su ecuación.


% =====================================================
\seccio{3}{Reglas de derivación}

\exf{1} Calcula la función derivada de estas funciones simples:\begin{tasks}(2)
  \task $f(x)=2x^3-5x^2+3x-2$
  \task $f(x)=\sqrt{2x}+\sqrt[3]{5x^2}$
  \task $f(x)=\dfrac{1}{x\sqrt x}$
  \task $f(x)=\dfrac{x^3}{x^2+1}$
  \task $f(x)=x^2\cdot\ln x$
  \task $f(x)=\sqrt{3x^3}\cdot e^x$
\end{tasks}\recorda{$(x^n)'=nx^{n-1}$ (también con exponentes fraccionarios: escribe las raíces como potencias). \ $(f\cdot g)'=f'g+fg'$; \ $\left(\frac fg\right)'=\frac{f'g-fg'}{g^2}$.}

\ex{1} Deriva y simplifica:\begin{tasks}(3)
  \task $f(x)=x^4-\dfrac{3}{x^2}+5$
  \task $f(x)=\dfrac{x+1}{x-1}$
  \task $f(x)=\dfrac{x^2}{x+2}$
  \task $f(x)=\dfrac{e^x}{x}$
  \task $f(x)=\log_2x+3^x$
  \task $f(x)=(2x-1)\,e^x$
  \task $f(x)=(x^2+1)\,e^x$
  \task $f(x)=\dfrac{\ln x}{x}$
  \task $f(x)=\dfrac{x}{\ln x}$
\end{tasks}\recorda{$(e^x)'=e^x$, \ $(a^x)'=a^x\ln a$, \ $(\ln x)'=\frac1x$, \ $(\log_ax)'=\frac{1}{x\ln a}$.}


% =====================================================
\seccio{4}{Regla de la cadena}

\exf{1} Calcula la función derivada de estas funciones compuestas:\begin{tasks}(2)
  \task $f(x)=\sqrt{x^2+1}$
  \task $f(x)=x\,e^{x^2}$
  \task $f(x)=\dfrac{2x}{(x^2+1)^2}$
  \task $f(x)=\ln(x^3-1)$
  \task $f(x)=\dfrac{x}{\sqrt{x^2+1}}$
  \task $f(x)=(2x+1)^4$
\end{tasks}\recorda{Regla de la cadena: $\left[f(g(x))\right]'=f'(g(x))\cdot g'(x)$. Derivamos la función «de fuera» dejando igual la «de dentro» y lo multiplicamos por la derivada de la de dentro.}

\ex{2} Deriva:\begin{tasks}(3)
  \task $f(x)=e^{3x^2-x}$
  \task $f(x)=(x^2-3x)^5$
  \task $f(x)=e^{\sqrt x}$
  \task $f(x)=\ln(e^x+1)$
  \task $f(x)=\dfrac{1}{(3x-2)^2}$
  \task $f(x)=\sqrt{\ln x}$
  \task $f(x)=\ln\dfrac{x-1}{x+1}$
  \task $f(x)=\log(x^2+1)$
  \task $f(x)=2^{x^2+1}$
\end{tasks}

\ex{2} Calcula $f'(1)$ si $f(x)=x^2e^{-x}$, y $g'(1)$ si $g(x)=\ln(x^2+3)$.


% =====================================================
\seccio{5}{Recta tangente y recta normal}

\exf{1} Calcula la ecuación de la recta tangente y de la recta normal a las funciones siguientes en el punto indicado:\begin{tasks}(1)
  \task $f(x)=x^2+2x-1$ \ en el punto de abscisa $x=1$
  \task $f(x)=2x^2-3x$ \ en el punto de abscisa $x=-1$
  \task $f(x)=\dfrac{x+1}{x-1}$ \ en el punto de abscisa $x=2$
\end{tasks}\recorda{Tangente: $y=f(a)+f'(a)(x-a)$. \ Normal (perpendicular a la tangente): $y=f(a)-\dfrac{1}{f'(a)}(x-a)$.}

\ex{1} Halla la recta tangente a $f(x)=\ln x$ en $x=1$ y la recta tangente a $g(x)=e^{2x}$ en $x=0$.

\ex{2} Dada $f(x)=x^3-3x^2+1$:\begin{tasks}(1)
  \task ¿En qué puntos la recta tangente es horizontal?
  \task ¿En qué puntos la recta tangente es paralela a la recta $y=9x$? Escribe la ecuación de estas tangentes.
\end{tasks}

\ex{2} Halla $a$ para que la recta tangente a $f(x)=x^2+ax$ en $x=1$ tenga pendiente 5. Escribe la ecuación de esta tangente.


% =====================================================
\seccio{6}{Monotonía y extremos}

\exf{1} Estudia la monotonía (intervalos de crecimiento y de decrecimiento) y los extremos relativos de las funciones:\begin{tasks}(2)
  \task $f(x)=-x^2+5x-1$
  \task $f(x)=2x^3-3x^2-12x+5$
\end{tasks}\recorda{Calculamos $f'(x)$, resolvemos $f'(x)=0$ y estudiamos el signo de $f'$ en los intervalos que determinan estos puntos y los que no son del dominio: donde $f'>0$ la función crece y donde $f'<0$ decrece. Si pasa de crecer a decrecer hay un máximo relativo; si pasa de decrecer a crecer, un mínimo.}

\ex{2} \label{ex:racional}Estudia la monotonía y los extremos relativos de:\begin{tasks}(3)
  \task $f(x)=\dfrac{x^2}{x-1}$
  \task $g(x)=x\,e^x$
  \task $h(x)=\ln(x^2+1)$
\end{tasks}

\ex{2} Halla los extremos relativos de $f(x)=x^4-8x^2+3$ con el criterio de la segunda derivada.

\ex{2} Calcula el máximo y el mínimo absolutos de $f(x)=x^3-3x+2$ en el intervalo $[-2,3]$.\recorda{Los extremos absolutos de una función continua en $[a,b]$ están en los puntos singulares de dentro del intervalo o en los extremos $a$ y $b$: basta con comparar sus valores.}

\ex{3} Halla $a$ y $b$ para que $f(x)=x^3+ax^2+b$ tenga un extremo relativo en el punto $(2,\,-3)$. ¿Es un máximo o un mínimo?


% =====================================================
\seccio{7}{Representación gráfica de funciones}

\exf{2} Estudia y representa la función $f(x)=x^3-3x$: dominio, puntos de corte, simetría, asíntotas, monotonía y extremos.\recorda{Pasos: dominio, puntos de corte, simetría, asíntotas, monotonía y puntos singulares, signo de la función, tabla de valores y representación gráfica.}

\ex{2} Representa la función $f(x)=\dfrac{x^2}{x-1}$ del ejercicio \ref{ex:racional}. Estudia antes sus asíntotas.

\ex{2} Estudia y representa $f(x)=\dfrac{4x}{x^2+1}$.

\ex{2} Estudia y representa $f(x)=x^4-2x^2$.

\ex{2} Estudia y representa $f(x)=x^3-6x^2+9x$.

\ex{2} Estudia y representa $f(x)=\dfrac{x^2+4}{x}$.

\ex{2} Estudia y representa $f(x)=\dfrac{x^2}{x^2-4}$.

\ex{3} Estudia y representa $f(x)=\dfrac{(x-1)^2}{x+1}$, con sus asíntotas verticales y oblicuas.

\ex{3} Estudia y representa $f(x)=\dfrac{x}{x^2-4}$. ¿Tiene extremos relativos?


% =====================================================
\seccio{8}{Problemas de optimización}

\exf{1} El beneficio mensual de una empresa, en miles de euros, cuando vende $x$ centenares de unidades es $B(x)=-x^2+40x-300$. Estudia el crecimiento y el decrecimiento de $B$. ¿Cuántas unidades tiene que vender para obtener el beneficio máximo? ¿Cuál es ese beneficio? ¿Para qué ventas tiene pérdidas?\recorda{Si conocemos la función: estudiamos el signo de la derivada (dónde crece y dónde decrece), hallamos los puntos donde $f'(x)=0$, justificamos si son máximos o mínimos y comprobamos que la solución tiene sentido en el contexto (unidades positivas, horas dentro del intervalo\ldots).}

\ex{2} El coste total de producir $x$ unidades de un artículo es $C(x)=x^2+100x+3\,600$ euros. Halla la función de coste medio por unidad, $M(x)=\frac{C(x)}{x}$, estudia su crecimiento y decrecimiento y calcula cuántas unidades hay que producir para que sea mínimo. ¿Cuál es ese coste medio?

\ex{2} El número de personas conectadas a una plataforma de vídeo (en miles) entre las 17 h y las 24 h es $N(t)=-2t^3+21t^2-36t+40$, donde $t\in[0,7]$ son las horas transcurridas desde las 17 h. Estudia cuándo aumenta y cuándo disminuye el número de conexiones. ¿A qué hora hay más personas conectadas? ¿Y menos?

\ex{2} Las ventas de un móvil nuevo, en miles de unidades al mes, $t$ meses después de lanzarlo son $V(t)=\dfrac{40t}{t^2+16}$, $t\geq0$. Estudia el crecimiento y el decrecimiento de las ventas. ¿En qué momento son máximas y cuántas se venden? ¿Qué pasa a largo plazo?

\ex{3} Una empresa estima que el beneficio de una campaña, en miles de euros, si invierte $x$ miles de euros en publicidad, es $B(x)=10x\,e^{-0{,}2x}$, $x\geq0$. Estudia el crecimiento y el decrecimiento de $B$. ¿Cuánto tiene que invertir para obtener el beneficio máximo? ¿Cuál es?

\ex{2} Una empresa vende cada unidad que fabrica a 120~€, y el coste de fabricar $x$ unidades al día es $C(x)=2x^2+40x+300$ euros. Escribe la función de beneficio diario y calcula cuántas unidades tiene que fabricar para maximizarlo. ¿Cuál es el beneficio máximo?\recorda{Si tenemos que construir la función: \ ingresos = precio $\cdot$ cantidad vendida; \ beneficio = ingresos $-$ costes. Elige bien la variable, escribe el precio y la cantidad en función de ella e indica qué valores puede tomar.}

\ex{2} Un hotel tiene 80 habitaciones y, a 50~€ la noche, las tiene todas ocupadas. Por cada euro que sube el precio, queda una habitación vacía. Cada habitación ocupada le cuesta 10~€ de limpieza y mantenimiento. ¿Qué precio hace máximo el beneficio? ¿Cuántas habitaciones se ocupan y cuál es el beneficio?

\ex{2} Un concierto benéfico tiene 800 espectadores si la entrada cuesta 20~€. Por cada euro que se rebaja la entrada, van 50 personas más. ¿Qué precio hace máxima la recaudación? ¿Cuántas personas van y cuánto se recauda?

\ex{3} La demanda semanal de un producto es $q=600-20p$ unidades cuando el precio es $p$ euros. Fabricar cada unidad cuesta 6~€ y, además, la empresa tiene unos costes fijos de 1\,000~€ semanales. Escribe el beneficio en función del precio y halla el precio que lo hace máximo. ¿Cuántas unidades se venden y cuál es el beneficio?

\ex{3} Una comerciante tiene 10\,000 kg de uva que podría vender hoy a 2~€ el kg. El precio aumentará 5 céntimos cada día, pero cada día se estropean 10 kg por cada cuadrado de los días que transcurren ($10t^2$ kg). Si puede esperar como mucho 20 días, ¿qué día le conviene vender?

\vfill
\begin{center}\color{muted}\small\faCheckCircle[regular]\ \ ¡Ánimo! Comprueba tus derivadas con la calculadora de los apuntes interactivos.\end{center}

\end{document}
